% Prose rendering of PrimeTensor/Structure.lean.
% Fragment: no \documentclass, no preamble, no document environment.

\section{Zero-free multiplicative carrier interface}
\label{Structure:sec}

\leanfile{PrimeTensor/Structure.lean}

\subsection*{What this file establishes}

By the end of \texttt{PrimeTensor/Algebra.lean} two carriers ---
$\name{MulRat}$ and its completion $\name{MulReal}$ --- satisfied the same six
laws, each proved twice, once on each side.  This file introduces the interface
that names those six laws, records both carriers as instances of it, and proves
the first theorem of the development that is stated once and holds for every
carrier at once.  It closes by checking that the embedding
$\name{ofRat} : \name{MulRat} \to \name{MulReal}$ respects all of the
structure.

The interesting content is a negative decision, and the module header states it
outright: Lean's $\ctor{Group}$ and $\ctor{CommGroup}$ are \emph{not} used, and
the class defined instead is deliberately weaker than they are.  Why a weaker
class is the right one here is the subject of
Design note~\ref{Structure:note:why-not-group}.

\subsection{The interface}

\begin{definition}[\lean{IsMulCarrier}]\label{Structure:def:ismulcarrier}
For a type $G$ already equipped with $\ctor{One}$, $\ctor{Mul}$ and
$\ctor{Inv}$ instances, $\name{IsMulCarrier}(G)$ is the proposition asserting
the six laws
\begin{align*}
  (ab)c &= a(bc), &
  1 \cdot a &= a, &
  a \cdot 1 &= a, \\
  ab &= ba, &
  (a^{-1})^{-1} &= a, &
  a \cdot a^{-1} &= 1 .
\end{align*}
\end{definition}

\begin{leancode}
class IsMulCarrier (G : Type*) [One G] [Mul G] [Inv G] : Prop where
  mul_assoc : ∀ a b c : G, (a * b) * c = a * (b * c)
  one_mul : ∀ a : G, 1 * a = a
  mul_one : ∀ a : G, a * 1 = a
  mul_comm : ∀ a b : G, a * b = b * a
  inv_inv : ∀ a : G, (a⁻¹)⁻¹ = a
  mul_inv : ∀ a : G, a * a⁻¹ = 1
\end{leancode}

Three features of the declaration carry the design.

It is a \emph{mixin}: the three operations appear as instance
\emph{parameters} in square brackets, not as fields.  A type does not acquire a
multiplication by being an $\name{IsMulCarrier}$; it must already have one, and
the class only asserts that the one it has obeys the laws.  So no competing
$\ctor{One}$, $\ctor{Mul}$ or $\ctor{Inv}$ instance is ever created, and there
is no possibility of two multiplications on the same carrier disagreeing about
which is meant.

It is \emph{$\ctor{Prop}$-valued}.  Every field is an equation, so an
$\name{IsMulCarrier}$ instance carries no data whatever --- nothing that could
be computed with, and nothing that could be unfolded to something unexpected.

And it is \emph{exactly six laws}, closed under nothing.  There is no division,
no power operation, no additive field, and no cancellation axiom.

\begin{designnote}[Why not $\ctor{CommGroup}$]\label{Structure:note:why-not-group}
The six laws say that $G$ is an abelian group, so $\ctor{CommGroup}\,G$ would
be provable and would bring the whole library with it.  The header rejects it
for the operations it would drag in.

Lean's $\ctor{Monoid}$ carries a natural-power field
$\ctor{npow} : \mathbb{N} \to M \to M$, and $\ctor{DivInvMonoid}$ carries an
integer power $\ctor{zpow} : \mathbb{Z} \to G \to G$ and a division
$\ctor{Div}$.  They are fields with defaults, not theorems: a $\ctor{CommGroup}$
instance on $\name{MulReal}$ would equip the carrier with $a^{0} = 1$, with
$a^{-n}$, and with $a / b$, whether or not any of them is ever written.

Each is precisely what this development has been at pains to exclude.
$\name{Depth}$ of \texttt{PrimeTensor/Depth.lean} has no zeroth constructor, so
that every axis is nonempty and every fold needs no identity element; a
$\ctor{npow}$ defined at $0$ reintroduces the zeroth index at the level of the
carrier and makes the exclusion cosmetic.  $\ctor{Div}$ reintroduces a second
primitive operation on top of $\ctor{Inv}$, so that a statement can be written
in two inequivalent-looking ways and every downstream $\ctor{simp}$ set has to
decide between them.  And an integer power presupposes $\mathbb{Z}$, an
additive object, in the object language of a development whose stated
constraint is that nothing additive appears.

The mixin is the smallest interface that gives the laws without the
operations.  Its cost is real and was visible in
\texttt{PrimeTensor/Scale/Laws.lean} and
\texttt{PrimeTensor/Algebra.lean}: since $\ctor{ac\_rfl}$, $\ctor{ring}$ and
$\ctor{group}$ all dispatch on the library's algebraic classes and not on this
one, associativity-commutativity regroupings still have to be written out by
hand, and both files did so.  The trade is spelled-out rewriting in exchange
for a carrier that carries no operation nobody asked for.
\end{designnote}

\begin{lemma}[\lean{IsMulCarrier.inv\_mul}]\label{Structure:lem:inv-mul}
In any $\name{IsMulCarrier}$, $a^{-1} \cdot a = 1$.
\end{lemma}

\begin{leancode}
theorem inv_mul (a : G) : a⁻¹ * a = 1 := by
  rw [IsMulCarrier.mul_comm]
  exact IsMulCarrier.mul_inv a
\end{leancode}

\begin{proof}
Commute and apply the sixth law.
\end{proof}

\begin{remark}[The first generic theorem]\label{Structure:rem:generic}
Trivial as it is, this is the first statement in the development proved once
for all carriers rather than once per carrier.  The same fact was proved
separately as $\name{MulRat.inv\_mul}$ in
\texttt{PrimeTensor/Scale/Laws.lean} and as $\name{MulReal.inv\_mul}$ in
\texttt{PrimeTensor/Algebra.lean}, each by this identical two-line argument.
Both survive: nothing is deleted, and the concrete names remain the ones the
earlier files' $\ctor{simp}$ sets refer to.  What changes is that a
\emph{third} carrier would now inherit the lemma instead of re-proving it.
\end{remark}

\begin{remark}[The axiom list is not minimal]\label{Structure:rem:redundant}
$\name{inv\_inv}$ follows from the other five.  Inverses in a monoid are
unique --- if $xy = 1$ and $xz = 1$ then
$y = y\cdot 1 = y(xz) = (yx)z = (xy)z = 1\cdot z = z$, using commutativity for
the fourth step --- and both $a$ and $(a^{-1})^{-1}$ are inverses of $a^{-1}$,
the first by Lemma~\ref{Structure:lem:inv-mul} and the second by the sixth law
applied to $a^{-1}$.  Listing it as an axiom costs one line per instance and
saves that derivation at every use, which for a class whose only two instances
already had the fact proved is the cheaper side of the trade.
\end{remark}

\subsection{The two instances}

\begin{proposition}[\lean{IsMulCarrier MulRat} and \lean{IsMulCarrier MulReal}]
\label{Structure:prop:instances}
Both $\name{MulRat}$ and $\name{MulReal}$ are $\name{IsMulCarrier}$s.
\end{proposition}

\begin{proof}
Each instance is a six-field record whose entries are the correspondingly named
theorems already proved: $\name{MulRat.mul\_assoc}$, $\name{one\_mul}$,
$\name{mul\_one}$, $\name{mul\_comm}$, $\name{inv\_inv}$, $\name{mul\_inv}$
from \texttt{PrimeTensor/MulRat.lean}, and the six theorems of the same names
on $\name{MulReal}$ from \texttt{PrimeTensor/Algebra.lean}.  Nothing is proved
here; the file only records that what was proved has the shape the interface
asks for.
\end{proof}

That both instances are literally lists of existing names is the evidence that
the interface was extracted from the practice rather than imposed on it: the
six laws are exactly the six that both carriers had already been given, in the
same order and with the same statements.

\subsection{The embedding respects the structure}

\begin{lemma}[\lean{MulReal.ofRat\_one}, \lean{ofRat\_mul} and \lean{ofRat\_inv}]
\label{Structure:lem:ofrat}
$\name{ofRat}(1) = 1$, and for all $a, b : \name{MulRat}$,
\[
  \name{ofRat}(a \cdot b) = \name{ofRat}(a) \cdot \name{ofRat}(b),
  \qquad
  \name{ofRat}(a^{-1}) = \name{ofRat}(a)^{-1} .
\]
\end{lemma}

\begin{leancode}
@[simp] theorem ofRat_mul (a b : MulRat) :
    ofRat (a * b) = ofRat a * ofRat b := by
  change
    Quotient.mk MulAsymptotic.setoid (MulCauchyStream.constant (a * b)) =
      Quotient.mk MulAsymptotic.setoid
        (MulCauchyStream.mul
          (MulCauchyStream.constant a)
          (MulCauchyStream.constant b))
  apply Quotient.sound
  apply MulAsymptotic.of_pointwise
  intro n
  rfl
\end{leancode}

\begin{proof}
The first is $\ctor{rfl}$: $\name{MulReal.one}$ was \emph{defined} as
$\name{ofRat}(1)$ in \texttt{PrimeTensor/Completion.lean} and installed as the
$\ctor{One}$ instance, so the two sides are the same term.

The other two are the four-line shape of
\texttt{PrimeTensor/Algebra.lean}: reduce the equality of classes to an
asymptotic equivalence by $\ctor{Quotient.sound}$, reduce that to a termwise
equality by $\name{of\_pointwise}$, and observe that the termwise equality is
$\ctor{rfl}$.  It is $\ctor{rfl}$ because the $n$-th term of the constant
stream at $ab$ is $ab$, and the $n$-th term of the product of the constant
streams at $a$ and at $b$ is also $ab$ --- the two streams are not merely
asymptotic, they are equal, and the same holds for inversion.
\end{proof}

\begin{remark}[Predicted, and for the predicted reason]\label{Structure:rem:predicted}
\texttt{PrimeTensor/Algebra.lean} closed by noting that the homomorphism
property of $\name{ofRat}$ was not yet stated but would follow from
$\name{of\_pointwise}$ in one line, since the constant stream of a product is
the product of the constant streams.  That is exactly the proof above, and the
``one line'' is the final $\ctor{rfl}$.  No analysis enters: the levels and
anchors of the completion play no part, because constant streams agree
identically rather than eventually.
\end{remark}

\begin{remark}[Injectivity is still missing]\label{Structure:rem:injective}
$\name{ofRat}$ is now known to preserve the pivot, the multiplication and the
inversion, which is everything the interface names.  It is still not known to
be injective, and injectivity is not a formal consequence of being a
homomorphism --- it is the statement that two rationals whose constant streams
are close at every level are equal, which is an archimedean fact about
$\name{MulRat}$ and needs an argument about the ladder.  Until it is proved,
$\name{MulReal}$ is a carrier receiving a structure-preserving map from
$\name{MulRat}$, not yet an extension of it.
\end{remark}

\begin{remark}[The mixin is not yet used generically]\label{Structure:rem:unused}
Lemma~\ref{Structure:lem:inv-mul} is the only theorem stated for a general
$\name{IsMulCarrier}$, and no later result in this file is.  The interface's
payoff is entirely deferred: it exists so that constructions downstream can be
written once against the laws instead of twice against the two carriers.  What
the file delivers now is the declaration, the two instances, and the evidence
that the embedding between them is a morphism of the structure the interface
describes.
\end{remark}

\subsection*{Summary of the correspondence}

\begin{center}
\small
\begin{tabular}{@{}lll@{}}
\hline
Lean declaration & Kind & Here \\
\hline
\lean{IsMulCarrier}                 & class      & Def.~\ref{Structure:def:ismulcarrier} \\
\lean{IsMulCarrier.inv\_mul}        & theorem    & Lem.~\ref{Structure:lem:inv-mul} \\
\lean{IsMulCarrier MulRat}          & instance   & Prop.~\ref{Structure:prop:instances} \\
\lean{IsMulCarrier MulReal}         & instance   & Prop.~\ref{Structure:prop:instances} \\
\lean{MulReal.ofRat\_one}           & simp thm   & Lem.~\ref{Structure:lem:ofrat} \\
\lean{MulReal.ofRat\_mul}           & simp thm   & Lem.~\ref{Structure:lem:ofrat} \\
\lean{MulReal.ofRat\_inv}           & simp thm   & Lem.~\ref{Structure:lem:ofrat} \\
\hline
\end{tabular}
\end{center}

\noindent
Every declaration of \texttt{PrimeTensor/Structure.lean} appears above,
together with the six fields of the class, listed in
Definition~\ref{Structure:def:ismulcarrier}.  The file contains no
\lean{sorry}, no axiom, and no unproved named proposition.
Remarks~\ref{Structure:rem:generic}, \ref{Structure:rem:redundant},
\ref{Structure:rem:predicted}, \ref{Structure:rem:injective}
and~\ref{Structure:rem:unused} are stated for the reader and are not
declarations of the file; the derivation in
Remark~\ref{Structure:rem:redundant} is not carried out in Lean.
